\section{Introduction}\label{sec_intro}

\subsection{Context and purpose}
%\bruno{Je ne pense pas que ce soit une super idée d'attaquer par le côté appliqué, ça lui donne trop de poids. C'est mieux d'attaquer sur le côté non-linéaire, et surtout adaptatif.} Pas Faux !

Time-frequency transforms and generalizations (wavelets and others) have long been used in various theoretical and applied domains. Besides quadratic transforms (Wigner distributions and generalizations), linear transforms such as the Gabor/STFT~\cite{daubechies1992ten,grochenig2013foundations} and wavelet transforms~\cite{Grossmann1984decomposition,daubechies1992ten,mallat2008wavelet,meyer1993wavelets} generally enjoy simple and useful invertibility properties, and therefore allow describing functions and signals as linear combination of building blocks, called time-frequency atoms. The time (and frequency/scale) resolution of the latter is specified by the construction rule: constant time and frequency resolution for Gabor/STFT, generated by translation and modulation, and constant relative frequency resolution for wavelets, generated by translation and scaling. See also~\cite{Kalisa1993weyl,Ali2000coherent,Fornasier2007banach} for alternative constructions that implement other scaling rules. Variants were also considered in specific applied domains, such as the Stockwell transform~\cite{Stockwell1996localization} in geophysics, which is very close to the constant $Q$ transform~\cite{Brown1991calculation,velasco2011constructing} we consider below, and the continuous wavelet transform.

In several application domains, in particular audio signal processing, it has been shown that adapting the scale of time-frequency atoms to the content of the signal can provide more efficient signal descriptions~\cite{Jaillet2007time,Liuni2013automatic,Leiber2023differentiable}. The window size is a main control parameter for the time-frequency resolution of the analysis: large windows provide good frequency resolution, while short windows yield good time resolution. The latter is constrained by the uncertainty principle, which can be given various quantitative formulations (see~\cite{Folland1997uncertainty,Ricaud2013refined} and references therein), and which basically states that precision in time domain is possible at the price of precision loss in frequency domain, and vice versa. The problem of tuning time-frequency resolution as a function of time or frequency has been addressed by various authors, in various contexts, often with pre-defined dependence, sometimes adaptively~\cite{Liuni2013automatic,Leiber2023differentiable}.  As a motivation, even though there is no complete consensus on psycho-physical aspects of human perception, it is known to involve several non-linear effects~\cite{Oxenham2018how}, and it has been claimed that this non-linearity allows going beyond time-frequency uncertainty in terms of localization~\cite{Oppenheim2013human}. Such strategy has been successfully implemented in some advanced audio coders such as AAC (see~\cite{Brandenburg1999mp3} for a short account), which can switch dynamically between short and long local cosine windows. Another possible motivation can be  hyper-resolution (separation of close locally harmonic components or close transients in signals for example). The adaptation is often driven by heuristic computations, for example the optimization of a sparsity criterion of the obtained time-frequency representation, using for example some form of entropy as in~\cite{Jaillet2007time,Liuni2013automatic} or other criteria~\cite{Leiber2023differentiable}. The adaptation may be implemented through a time (or frequency) dependent warping function, as in~\cite{Holighaus2019continuous}. To our knowledge, the non-linear problem where the time-frequency resolution is adapted to the analyzed function, has not been analyzed from the mathematical point of view so far.

\smallskip
The goal of the present paper is to introduce and study such adaptive time-frequency transforms, able to adapt their time-frequency resolution to the analyzed function.
This is done here by introducing a \textit{focus function} $f\mapsto\sigma_f$ which adapts the shape (size, bandwidth etc...) of the analysis window to specific properties of the analyzed signal $f$. %, and then then allows the transform to adapt themselves to the signal. 
Stepping away from fixed time-frequency resolution makes the analysis significantly more complex.
The purpose of this article is to introduce the non-linear transforms, prove that they are well-defined on $L^2$ and provide explicit, signal-dependent, lower and upper bounds for their norm (which depends on the choice of focus function).

We first introduce in Section~\ref{sec_time_focus} a time-focused transform $M^\tau$, which is a modified STFT where the window scale can be adapted to the signal at each time. The adaptation is done by associating with the analyzed signal $f$ a focus function $\sigma_f^\tau:\,t\mapsto \sigma^\tau_f(t)$. We prove in Theorem~\ref{theo_bounding_time_marseillan_transform} the well-definedness of $M^\tau$ as a map from $L^2(\R)$ into $L^2(\R^2)$ and obtain a norm control of the form $c_f\|f\|^2\le \|M^\tau f\|^2\le C_f\|f\|^2$, with explicit constants $c_f,C_f$, under suitable assumptions on the focus function $\sigma_f^\tau$. Building a frequency-focused STFT can be done along similar lines, and is not addressed here.
After briefly reviewing constant-Q transform~\cite{velasco2011constructing,holighaus2012framework} and continuous wavelet transform~\cite{Grossmann1984decomposition} and pointing out their relationship, we introduce in Section~\ref{sec_frequency_focus} a new frequency-focused transform $M^\nu$. The latter is built from a complex-valued wavelet transform, modified by a focus function defined in the frequency domain: to the analyzed signal $f$ is associated a focus function $\sigma_f^\nu: \,\omega\mapsto\sigma^\nu_f(\omega)$. We prove similar norm control and well-definedness results.
%Then in Section~\ref{sec_CSQT} we detail a generalisation of a constant-Q transform studied in~\cite{velasco2011constructing,holighaus2012framework}, and prove an isometry result, using notations and ideas that will be useful for the rest of the paper. }
Explicit examples of focus functions are discussed in the context of continuous time, in Section~\ref{sec_time_focus} for time focus, which can be transposed to the frequency focus case. Section~\ref{se:numerical.experiments} is devoted to numerical illustrations, in discrete, finite-dimensional situations. There, we discuss and display examples of focus functions and illustrate the resulting time-focused and frequency-focused spectrograms, computed on real audio signals and toy examples.
Section~\ref{sec_conclusion} is devoted to conclusions and perspectives.

\smallskip
Except the section devoted to numerical illustrations, the analysis described in the present paper is limited to the continuous time setting and mostly $L^2(\R)$. Extensions to more general functional settings and issues related to discretization of the transforms will be the object of further work.





\subsection{Notation}

We first introduce or recall some notation. We will often use the notation $C([a,b],X)$ for the space of continuous functions from $\R$ into $X$, supported  in $[a,b]$. More generally, given $A\subset \R^\R$ we denote by $A_c$ the subspace of functions in $A$ with compact support, and by $A_0$ the subspace of functions in $A$ which vanish at infinity. We define $A+c:=\{f+c, f\in A\}$ for $c\in\R$ and $A^+$ denotes the subspace of $A$ of non-negative valued functions.

Given an open interval $I$ in $\R$, $C^n(I)$ stands for the space of function which are $n$ times continuously differentiable on $I$, and $C^\infty(I) = \cap_{n\geqslant 0} C^n(I)$. Furthermore, we denote  by $C_p(I)$ the space of piecewise continuous functions on $I$, and $C_{p,0}(I)$ the subspace of $C_p(I)$ consisting of functions that tend to 0 at $\pm\infty$. $C_c(I)$ denotes the space of continuous, compactly supported functions on $I$.
%Finally, we denote by $\mathcal{E}$ the space of step functions, and by $\mathcal{E}([a,b])$ the space of step function with support included in $[a,b]$.

We recall that $L^p(\R^d,d\mu)$, $p\geqslant1$, stands for the set of $p-$integrable functions - where we identify functions that coincide almost everywhere - with respect to the measure $\mu$. The shorter notations $L^p(\R^d)$ and $L^p$ denote respectively $L^p(\R^d,dx)$ and $L^p(\R,dx)$ where $dx$ is the Lebesgue measure.
%When several variable are involved the notation $L^p_x$ is used to specify that the $L^p$ norm is taken with respect to the $x$ variable. When several $L^p$ norm will be involved for different variable then the following mixed norm notation will be used, for some $p,q\in[1,+\infty]$
%\begin{equation}
% \|f\|_{L^p_x L^q_t}:= \|\|f(t,x)\|_{L^q_t}\|_{L^p_x}\ .
%\end{equation}


Given $f\in L^p(\R^d)$, its Fourier transform is either written $\F(f)$ or $\hat{f}$, and defined by the following convention.  For $f\in L^1(\R^d)$,
\begin{equation}
%    \forall \omega\in\R^d,\quad 
    \F(f)(\omega) := \int_{\R^d} f(x)e^{-2i\pi x\cdot\omega} dx\ ,\qquad\omega\in\R^d\ .
\end{equation}
With this definition, the inverse Fourier transform reads
\begin{equation}
\mathcal{F}^{-1}(F)(t):= \int_{\R^d}F(\omega)e^{2i\pi t\cdot \omega}d\omega\ , \qquad t\in\R^d\ .
\end{equation}


For functions of two variables, we introduce the notation
\begin{align}
   \forall x_1,x_2,\xi_1,\xi_2\in\R, \qquad \F_1 (f)(\xi_1,x_2)&:= \int_\R f(x_1,x_2)e^{-2i\pi x_1\xi_1}dx_1 \ ; \\
   \F_2 (f)(x_1,\xi_2)&:= \int_\R f(x_1,x_2)e^{-2i\pi x_2\xi_2}dx_2\ ,
\end{align}
and
\begin{align}
    \forall  x_1,x_2,\xi_1,\xi_2\in\R, \qquad \F^{-1}_1 (F)(x_1,\xi_2)&:= \int_\R F(\xi_1,\xi_2)e^{2i\pi x_1\xi_1}d\xi_1 \ ; \\
    \F^{-1}_2 (F)(\xi_1,x_2) & := \int_\R F(\xi_1,\xi_2)e^{2i\pi x_2\xi_2}d\xi_2\ .
\end{align}
